\bookchapter{Bab 18: Pencadangan dan migrasi data}

\chapterintro{1. Pertama pahami: Apa saja di ponsel Anda yang hilang akan hilang selamanya}

\booksection{“Status kelangsungan” data yang berbeda tidak sama, Anda harus membedakannya:}

Tersimpan lokal di ponsel, hilang berarti benar-benar hilang

\begin{itemize}
  \item \textbf{: Foto video (jika tidak mengaktifkan sinkronisasi cloud), catatan, SMS, riwayat panggilan, file yang diunduh.}Tersimpan di server aplikasi, ganti ponsel login ulang akan ada
  \item \textbf{: Daftar teman WeChat, pesanan Taobao, riwayat transaksi di aplikasi bank. Hal-hal ini disimpan di server perusahaan terkait, bukan di ponsel Anda.}Di zona abu-abu, perlu operasi khusus
  \item \textbf{: Riwayat chat WeChat (meskipun server WeChat menyimpannya, tetapi Anda sendiri tidak dapat mengunduh dari server), kredensial perangkat yang terikat ke mobile banking, kunci di aplikasi autentikasi dua langkah (seperti Google Authenticator).}Jadi yang perlu Anda lindungi dengan prioritas adalah kategori pertama dan ketiga.
\end{itemize}

2. Pencadangan seluruh perangkat: Selesai semua dalam sekali

\booksection{Pencadangan seluruh perangkat seperti mengambil “snapshot” lengkap ponsel — pengaturan sistem saat ini, tata letak desktop, data aplikasi, foto, kontak, semuanya dikemas.}

Opsi 1: Pencadangan lokal komputer (paling andal, direkomendasikan untuk yang punya komputer)

\booksubsection{iPhone}

\begin{itemize}
  \item \textbf{: Pengaturan Nama Anda iCloud iCloud Cloud Backup Aktifkan. Ponsel terhubung daya, terhubung Wi-Fi, saat layar terkunci akan otomatis dicadangkan. 5GB gratis kemungkinan besar tidak cukup, disarankan upgrade ke 50GB dengan 6 yuan.}Android
  \item \textbf{Android}Keuntungan pencadangan komputer adalah tidak perlu membeli ruang cloud, dan kecepatan pemulihan jauh lebih cepat dari pemulihan cloud. Kerugiannya adalah Anda harus ingat untuk secara berkala mencolokkan kabel untuk mencadangkan — sebagian besar orang tidak bisa melakukan ini.
  \item Opsi 2: Pencadangan cloud (paling nyaman, direkomendasikan untuk semua orang)
\end{itemize}

\booksubsection{iPhone}

\begin{itemize}
  \item \textbf{: Pengaturan Nama Anda iCloud iCloud Cloud Backup Aktifkan. Ponsel terhubung daya, terhubung Wi-Fi, saat layar terkunci akan otomatis dicadangkan. 5GB gratis kemungkinan besar tidak cukup, disarankan upgrade ke 50GB dengan 6 yuan.}Android
  \item \textbf{Android}Keuntungan terbesar pencadangan cloud adalah "tidak perlu Anda pusing", otomatis selesai. Kerugian terbesar adalah saat pemulihan harus mengunduh lama, dan ruang gratis umumnya tidak cukup.
  \item 3. Strategi pencadangan terkategori: Hal penting disimpan terpisah
\end{itemize}

\booksection{Pencadangan seluruh perangkat adalah jaring pengaman, tetapi beberapa hal terlalu penting, layak disimpan terpisah:}

Foto dan video (aset digital paling berharga)

\booksubsection{Skema ganda}

\begin{itemize}
  \item \textbf{: Aktifkan sinkronisasi otomatis iCloud Foto/Google Foto, sambil secara berkala setiap beberapa bulan menyalin foto ke komputer atau hard disk eksternal. Jadi satu di cloud, satu di lokal, saling mencadangkan.}Skema tanpa biaya
  \item \textbf{: Google Foto memberikan 15GB gratis, dapat diatur ke "kualitas terkompresi" untuk menghemat ruang. Atau setiap bulan gunakan kabel data untuk menyalin semua foto ke komputer.}Kontak (mudah diabaikan tetapi sangat penting)
\end{itemize}

\booksubsection{Kontak iPhone secara default disinkronkan dengan iCloud, selama sinkronisasi aktif tidak perlu pencadangan ekstra. Tetapi jika ingin lebih aman, bisa ekspor file vCard di iCloud.com simpan di komputer.}

\begin{itemize}
  \item Kontak Android biasanya disinkronkan dengan Google atau cloud merek. Ekspor file .vcf di pengaturan aplikasi Kontak, itu adalah cadangan lengkap.
  \item Riwayat chat WeChat (ditangani terpisah, yang paling penting)
\end{itemize}

\booksubsection{Ini harus dibahas terpisah, karena dalam pencadangan seluruh perangkat riwayat chat WeChat sering tidak lengkap, dan alat pindah ponsel sama sekali tidak bisa memindahkannya.}

Metode 1: Migrasi ke ponsel baru

\begin{itemize}
  \item \textbf{. WeChat Saya Pengaturan Umum Migrasi \& Pencadangan Riwayat Chat Migrasi Migrasi ke Ponsel/Tablet. Ponsel lama dan baru terhubung Wi-Fi yang sama, pindai kode untuk transfer. Catatan: Sebelum migrasi selesai, ponsel lama jangan keluar WeChat.}Metode 2: Cadangkan ke komputer
  \item \textbf{. Pasang WeChat versi PC di komputer, tiga garis di kiri bawah Migrasi \& Pencadangan Pencadangan \& Pemulihan Cadangkan riwayat chat ke komputer. Setelah itu ganti ponsel bisa dipulihkan ke ponsel baru.}Metode 3: Cadangkan ke cloud
  \item \textbf{. WeChat juga mendukung pencadangan ke iCloud atau ruang cloud ponsel, tetapi perhatikan: cadangan ini hanya bisa dipulihkan di ponsel saat ini, tidak bisa dipulihkan lintas perangkat. Untuk ganti ponsel, masih harus menggunakan dua metode sebelumnya.}4. Panduan ganti ponsel: Dari ponsel lama ke ponsel baru
\end{itemize}

\booksection{Ganti ponsel baru adalah hal yang menyenangkan, tetapi jika data tidak dipindahkan menjadi sakit kepala. Tiga situasi ditangani secara terpisah:}

Upgrade merek yang sama (paling sederhana)

\booksubsection{iPhone iPhone baru}

\begin{itemize}
  \item \textbf{: Saat menyalakan aktivasi letakkan ponsel lama di sebelahnya, akan otomatis muncul "Quick Start", gunakan ponsel lama untuk memindai pola di ponsel baru. Semua data, pengaturan, tata letak aplikasi akan dipindahkan apa adanya, termasuk kata sandi Wi-Fi akan otomatis terbawa. Membutuhkan sekitar 20-30 menit.}Android merek yang sama
  \item \textbf{: Gunakan alat pindah ponsel bawaan merek. Huawei gunakan "Phone Clone", Xiaomi gunakan "Xiaomi Mover", OPPO gunakan "Phone Clone", vivo gunakan "EasyShare". Langkahnya hampir sama: ponsel lama dan baru pasang aplikasi yang sama, pindai kode untuk terhubung, centang yang ingin dipindahkan, mulai transfer.}Ganti ponsel lintas merek (Android ke Android, merek berbeda)
\end{itemize}

\booksubsection{Gunakan alat pindah ponsel pihak ketiga atau gunakan aplikasi pindah ponsel bawaan merek lama (aplikasi pindah ponsel sebagian besar merek juga mendukung pindah ke ponsel Android lain). Yang bisa dipindahkan: kontak, foto, SMS, riwayat panggilan, sebagian aplikasi. Yang tidak bisa dipindahkan: status login dalam aplikasi (Anda harus login ulang), riwayat chat WeChat (gunakan fungsi migrasi WeChat sendiri).}

Ganti ponsel lintas platform (Android ke iPhone, atau sebaliknya)

\booksubsection{Android iPhone}

\begin{itemize}
  \item \textbf{: Apple punya aplikasi "Move to iOS", unduh di ponsel Android, saat aktivasi iPhone baru pilih "Pindahkan data dari perangkat Android". Bisa memindahkan kontak, foto, SMS dan data dasar lainnya, tetapi data aplikasi tidak bisa dipindahkan. Riwayat chat WeChat masih harus menggunakan fungsi migrasi WeChat sendiri.}iPhone Android
  \item \textbf{: Setiap merek Android memiliki alat migrasi yang sesuai, tetapi sejujurnya, pengalaman migrasi lintas platform tidak sebaik platform yang sama. Hal paling penting dilakukan dulu secara terpisah: ekspor kontak, cadangkan foto, migrasi riwayat chat WeChat. Sisanya perlahan-lahan.}5. Aturan emas pencadangan: Prinsip 3-2-1
\end{itemize}

\booksection{Ini adalah prinsip klasik di bidang keamanan data, sederhana dan mudah diingat:}

3

\begin{itemize}
  \item \textbf{: Data penting simpan setidaknya tiga salinan (satu asli, dua cadangan).}2
  \item \textbf{: Cadangan gunakan dua media penyimpanan yang berbeda (misalnya satu di cloud, satu di hard disk eksternal).}1
  \item \textbf{: Setidaknya satu cadangan diletakkan di lokasi fisik yang berbeda (mencegah kebakaran, pencurian, dll. menghancurkan semua data sekaligus).}Untuk orang biasa, mencapai "satu di cloud, satu di komputer" sudah dianggap sangat baik. Tidak perlu terlalu rumit, yang penting Anda benar-benar melakukannya.
\end{itemize}

Perhatian khusus riwayat chat WeChat

\begin{tcbwarn}{Riwayat chat WeChat tidak dapat dimigrasi melalui alat pindah ponsel apa pun! Itu tidak termasuk dalam cakupan pencadangan seluruh perangkat. Harus menggunakan fungsi "Migrasi Riwayat Chat" bawaan WeChat, atau cadangkan ke WeChat versi PC. Banyak orang setelah pindah ke ponsel baru, baru menemukan riwayat chat semua hilang, saat itu sudah terlambat — ponsel lama mungkin sudah di-reset pabrik.}
Berlatih: Lakukan latihan pencadangan lengkap

\end{tcbwarn}

\begin{tcbpractice}{Periksa apakah iCloud Cloud Backup atau pencadangan cloud merek aktif, konfirmasikan waktu pencadangan terakhir.}
\begin{enumerate}
  \item Cadangkan foto secara manual: gunakan kabel data hubungkan ke komputer, salin folder DCIM dan Pictures ke komputer.
  \item Cadangkan kontak: pengguna iPhone ekspor vCard di iCloud.com; pengguna Android ekspor file .vcf di pengaturan aplikasi Kontak.
  \item Cadangkan riwayat chat WeChat: buka WeChat Saya Pengaturan Umum Migrasi \& Pencadangan Riwayat Chat Cadangkan ke komputer.
  \item Salin file vcf kontak dan folder foto yang diekspor ke hard disk eksternal atau USB flash drive, menyelesaikan "satu offline di tempat lain" dari prinsip "3-2-1".
  \item Bab 19: Pemecahan masalah dan penyelamatan diri
\end{enumerate}

\end{tcbpractice}
